\scanpage{61}
19) Considérons \deleted{\unsure{$G$ un}} \inserted{graphe} cubique $\Delta$ avec partition \unsure{trialitaire|trilétaire|tritétraèdre} $S=\bigcup_{\alpha\in\Lambda}X_\alpha$, et considérons le \unsure{gpe|groupe} $G\ (\simeq\unsure{\mathfrak S_3|\mathscr S_3}\mathbin{.}(\mathfrak S_3\times\mathfrak S_3\times\mathfrak S_3))$ des autom de $S$ ainsi décomposé, $\deleted{\unsure{G}}$ et son opération sur l'ens \inserted{$(\mathrm{Tria})$} des triades de $\Delta$, \inserted{et sur l'ens des (\unsure{bitr|bitri}) des bitriangles de $\Delta$.} Alors \deleted{il y a}

a) \deleted{existent} \inserted{Il y a} \deleted{trois} \inserted{quatre} orbites de $G$ sur Tria :

1°) L'orbite formée des triades contenues dans un des $X_\alpha$

2°) L'orbite \deleted{formée} des triades de la forme
\diagram{\begin{tikzcd}[ampersand replacement=\&,row sep=small,column sep=large]
\bullet\arrow[r,dash]\&\bullet\arrow[ddl,dash]\&\bullet\&\\
\bullet\&\bullet\arrow[r,dash]\&\bullet\&\text{(type I)}\\
\bullet\&\bullet\&\bullet\&
\end{tikzcd}}

3°) L'orbite \deleted{formée} des triades de la forme
\diagram{\begin{tikzcd}[ampersand replacement=\&,row sep=small,column sep=large]
\bullet\arrow[r,dash]\arrow[dr,dash]\&\bullet\&\bullet\&\\
\bullet\&\bullet\&\bullet\&\text{(type II)}\\
\bullet\&\bullet\arrow[r,dash]\&\bullet\&
\end{tikzcd}}

4°) L'orbite des triades de la forme
\diagram{\begin{tikzcd}[ampersand replacement=\&,row sep=tiny,column sep=small]
\bullet\arrow[rr,dash,bend left=12]\&\bullet\&\bullet\&\\
\bullet\&\bullet\arrow[r,dash]\&\bullet\&\text{(type III)}\\
\bullet\arrow[r,dash]\&\bullet\&\bullet\&
\end{tikzcd}}

b) Les triades contigues à une triade de type II sont de type I, celles contigues à une triade de type I sont de type I l'une, de type II l'autre de sorte que \deleted{\unsure{un}} l'hexagone \deleted{de triades} \inserted{d'une triade de type I ou II} est de type
\diagram{\begin{tikzcd}[ampersand replacement=\&,row sep=tiny,column sep=small]
\&\mathrm I\arrow[r,dash]\&\mathrm{II}\arrow[dr,dash]\&\\
\mathrm{II}\arrow[ur,dash]\arrow[dr,dash]\&\&\&\mathrm{II}\arrow[dl,dash]\\
\&\mathrm{II}\arrow[r,dash]\&\mathrm I\&
\end{tikzcd}}
(2 sommets \inserted{opposés} de type I, les autres de type II), l'hexagone d'une triade de type (III) est \inserted{formé de triades} de type III

% archive sheet number: 60

